\chapter{Coût, mesure de l'erreur}
\label{ch:erreur}

\begin{lecture}
Déplacer la frontière par le calcul suppose de savoir si elle est mal placée, et de
combien. Un seul nombre doit résumer les mille six cents prédictions. Ce nombre
est le coût, et il ne dépend que d'une quantité : la probabilité que le modèle
accordait à la réponse observée.
\end{lecture}

\section{Coût d'un élève}

Le fichier donne l'élève $3$ pour \Gryf, et le modèle lui accordait
$0{,}79$. Cette probabilité est la quantité à évaluer : plus elle est proche de
$1$, plus la prédiction est correcte.

L'élève $1$ est \Slyth, et le modèle annonçait $0{,}016$ de probabilité
d'être \Gryf, donc $0{,}984$ de ne pas l'être. C'est cette seconde valeur
qui entre dans le calcul.

Le coût d'un élève est l'opposé du logarithme de cette probabilité. La fonction ainsi obtenue s'appelle l'\emph{entropie croisée binaire} : elle
croise deux distributions sur les deux réponses possibles, celle qu'observe le
fichier et celle qu'annonce le modèle, et mesure ce que la seconde coûte lorsque
la première est vraie. Son autre nom, log-vraisemblance négative, vient de sa
construction (ch.~\ref{ch:cout}).
\[
  \ell=-\log(\text{probabilité accordée à la bonne réponse}).
\]
Le logarithme transforme un produit en somme, ce qui est nécessaire pour agréger
mille six cents élèves : le produit d'autant de facteurs inférieurs à $1$ passe
sous le plus petit flottant représentable.

\begin{table}[htbp]
\centering
\begin{tabular}{llrrr}
\toprule
élève & maison réelle & $p$ (\Gryf) & prob.\ de la vérité & coût \\
\midrule
$3$ & \Gryf & $0{,}792$ & $0{,}792$ & $0{,}233$ \\
$4$ & \Gryf & $0{,}991$ & $0{,}991$ & $0{,}009$ \\
$1$ & \Slyth  & $0{,}016$ & $0{,}984$ & $0{,}016$ \\
$0$ & \Rav  & $0{,}001$ & $0{,}999$ & $0{,}001$ \\
\bottomrule
\end{tabular}
\caption{Les quatre élèves du chapitre précédent. Aucun n'est mal classé, donc
tous les coûts sont faibles ; le plus élevé est celui de l'élève $3$, le cas
limite.}
\label{tab:cout-eleves}
\end{table}

Une prédiction correcte assortie d'une probabilité élevée coûte presque zéro. Une prédiction fausse assortie de la
même probabilité coûte d'autant plus que celle-ci est élevée, et sans borne
supérieure : le modèle ne peut donc pas maintenir une prédiction fausse avec une
forte probabilité.

\begin{figure}[htbp]
\centering
\begin{tikzpicture}
\begin{axis}[courseplot,height=6cm,axis lines=left,
  xlabel={probabilité accordée à la bonne réponse},ylabel={coût},
  xmin=0,xmax=1,ymin=0,ymax=5,samples=250,domain=.008:1]
  \addplot[accent,very thick]{-ln(x)};
  \addplot[only marks,mark=*,warning,mark size=2.4pt]
    coordinates {(0.792,0.233) (0.5,0.693) (0.1,2.303)};
  \node[warning,anchor=south west,font=\small] at (axis cs:.5,.78){$\log 2=0{,}693$};
  \node[warning,anchor=west,font=\small] at (axis cs:.1,2.5){$p=0{,}1$};
\end{axis}
\end{tikzpicture}
\caption{Le coût d'un seul élève. Il vaut $\log 2$ en $0{,}5$, tend vers zéro à
droite et croît sans borne à gauche.}
\label{fig:cout-un}
\end{figure}
\FloatBarrier

\section{Coût moyen sur le fichier}

Le coût précédent juge une prédiction. Placer la frontière demande de juger un
modèle, c'est-à-dire les mille six cents prédictions qu'il produit, et de le
comparer à un autre modèle. Mille six cents coûts individuels ne se comparent pas : déplacer la frontière
améliore en effet les uns et dégrade les autres, et deux modèles se retrouvent
meilleurs chacun sur une partie du fichier. Le calcul a
besoin d'une quantité unique, qui tranche.

\begin{figure}[htbp]
\centering
\begin{tikzpicture}[font=\footnotesize,
  row/.style={anchor=base,font=\footnotesize},
  val/.style={anchor=base east,font=\footnotesize}]
\node[row,font=\scriptsize\bfseries] at (0,0.75) {élève};
\node[val,font=\scriptsize\bfseries] at (2.5,0.75) {$y_i$};
\node[val,font=\scriptsize\bfseries] at (4.0,0.75) {$p_i$};
\node[val,font=\scriptsize\bfseries] at (6.0,0.75) {coût $\ell_i$};
\draw[ink!30] (-0.3,0.55) -- (6.2,0.55);
\foreach \k/\e/\y/\p/\c in {0/3/1/0{,}792/0{,}233, 1/4/1/0{,}991/0{,}009,
                            2/1/0/0{,}016/0{,}016, 3/0/0/0{,}001/0{,}001}{
  \node[row] at (0,-0.52*\k) {$\e$};
  \node[val] at (2.5,-0.52*\k) {$\y$};
  \node[val] at (4.0,-0.52*\k) {$\p$};
  \node[val] at (6.0,-0.52*\k) {$\c$};
}
\node[row,text=ink!60] at (0,-2.14) {$\vdots$};
\node[val,text=ink!60] at (6.0,-2.14) {$\vdots$};
\node[row,text=ink!70,font=\scriptsize] at (0.35,-2.66) {$1\,600$ lignes};
\draw[decorate,decoration={brace,amplitude=5pt},ink!45]
  (6.45,0.35) -- (6.45,-2.8);
\node[anchor=west,font=\small] at (7.0,-1.2)
  {$J=\dfrac{1}{m}\displaystyle\sum_{i=1}^{m}\ell_i=0{,}108$};
\end{tikzpicture}
\caption{Les quatre élèves du tableau~\ref{tab:cout-eleves}, suivis des mille cinq
cent quatre-vingt-seize autres. Chacun reçoit son propre coût, et la moyenne les
réduit à un seul nombre : celui que la descente fait baisser, et le seul qui
permette de comparer deux jeux de coefficients.}
\label{fig:agregation}
\end{figure}
\FloatBarrier

La moyenne remplit cet office, et le fait mieux que la somme : celle-ci croît
avec le nombre d'élèves, si bien que le même modèle afficherait un coût dix fois
plus grand sur un fichier dix fois plus long. Divisée par $m$, la valeur reste
comparable d'un fichier à l'autre --- entre l'entraînement et le test, notamment
--- et la longueur du pas de la descente conserve le même sens quel que soit le
nombre d'élèves.

En notant $y_i$ la réponse, valant $1$ pour un \Gryf et $0$ sinon, et $p_i$ la
probabilité annoncée :
\[
  J=-\ann{rulecolor}{\frac{1}{m}\sum_{i=1}^{m}}{\annl{moyenne sur}\\ \annl{les élèves}}
    \Bigl[\ann{accent}{y_i\log p_i}{\annl{seul terme retenu}\\ \annl{si $y_i=1$}}
    +\ann{warning}{(1-y_i)\log(1-p_i)}{\annl{seul terme retenu}\\ \annl{si $y_i=0$}}\Bigr]
\]

Un seul des deux termes agit pour chaque élève. Pour un \Gryf, $y_i=1$ : le
second s'annule et il reste $-\log p_i$. Pour un élève d'une autre maison,
$y_i=0$ : le premier s'annule et il reste $-\log(1-p_i)$. Le facteur $y_i$ joue ainsi le rôle d'un interrupteur, et cette écriture en une seule ligne remplace un
branchement, ce qui la rend dérivable d'un bout à l'autre --- propriété dont le
chapitre suivant a besoin.

\begin{resultat}[Coût au point de départ]
Si tous les coefficients valent zéro, tous les scores valent zéro, donc toutes
les probabilités valent $0{,}5$. Chaque élève coûte alors $-\log(0{,}5)$, et la
moyenne aussi :
\[
  J=\log 2=0{,}693147 .
\]
C'est la première valeur qu'affiche n'importe quel entraînement partant de zéro.
Une autre valeur situe l'erreur dans le calcul du score ou dans
l'initialisation.
\end{resultat}

Reste à dire ce qu'on fait de ce nombre. Il dépend des coefficients, seule chose
qui varie une fois le fichier fixé : les probabilités $p_i$ viennent des scores,
les scores viennent des coefficients. Déplacer la frontière change les mille six cents probabilités, partant les mille
six cents coûts, et avec eux leur moyenne. Le coût
s'écrit ainsi comme une fonction des trois coefficients, et se note $J(w)$.

Cette dépendance transforme la question de départ. Placer une frontière dans un
nuage de points est un problème de géométrie, sans procédé de résolution évident.
Chercher les trois nombres qui rendent $J(w)$ le plus petit possible est un
problème de minimisation d'une fonction de trois variables, pour lequel il existe
des méthodes. Les deux questions ont la même réponse, et c'est la seconde que le
programme résout.

\begin{resultat}[Formulation du problème]
\[
  \underbrace{\text{où placer la frontière ?}}_{\text{géométrie}}
  \qquad\longrightarrow\qquad
  \underbrace{\text{quel } w \text{ minimise } J(w)\text{ ?}}_{\text{calcul}}
\]
\end{resultat}

La surface que définit $J(w)$ au-dessus des couples de coefficients possibles
apparaît à la figure~\ref{fig:surface} ; la méthode consiste à en suivre la pente
vers le bas.

Deux questions restent ouvertes sur cette moyenne, et les deux sections suivantes
les traitent : ce qu'elle recouvre, mille six cents coûts pouvant se répartir
très différemment pour une même valeur, et ce qui la rend préférable au comptage
des élèves bien classés.

\section{Répartition du coût}

Le coût moyen final vaut $0{,}108$, et cette valeur agrège mille six cents
nombres très inégaux entre eux. L'inégalité tient à la forme du coût individuel.
Un élève bien classé et loin de la frontière reçoit une probabilité proche de
$1$, dont le logarithme est proche de zéro : sa contribution s'éteint. Un élève
du mauvais côté reçoit une probabilité d'autant plus petite qu'il en est loin, et
$-\log p$ croît alors sans borne. La moyenne se trouve portée par les quelques élèves que le modèle place du mauvais côté avec assurance.

\begin{figure}[htbp]
\centering
\begin{tikzpicture}
\begin{axis}[cloud,height=8cm,
  xlabel={\Herb, en écarts-types},ylabel={\Runes, en écarts-types},
  xmin=-2.8,xmax=2.8,ymin=-2.8,ymax=2.8,
  colorbar,colorbar style={ylabel={coût de l'élève},font=\small},
  colormap={cout}{color=(gridgray!45) color=(warning!60) color=(warning)},
  point meta min=0,point meta max=3,legend style={draw=none}]
  \addplot[scatter,only marks,mark size=1.5pt,scatter src=explicit,opacity=.85]
    table[x=x,y=y,meta=cost] {data/cout_eleves.dat};
  \addplot[ink,very thick,domain=-2.8:2.8]{(4.7454+3.4535*x)/3.4688};
\end{axis}
\end{tikzpicture}
\caption{Chaque élève, coloré par son coût dans le modèle final. La quasi-totalité
du nuage est grise, ces élèves contribuant de façon négligeable. Seuls quelques
points, tous situés du côté opposé à leur maison, portent une couleur.}
\label{fig:cout-nuage}
\end{figure}
\FloatBarrier

\begin{chiffres}[Concentration du coût]
\begin{center}
\begin{tabular}{lr}
\toprule
part du coût portée par les $2\,\%$ les plus coûteux & $60{,}5\,\%$ \\
coût de l'élève le plus coûteux & $11{,}1$ \\
coût moyen & $0{,}108$ \\
\bottomrule
\end{tabular}
\end{center}
Cet élève est un \Gryf dont la note d'\Herb dépasse la moyenne et dont celle
d'\Runes lui est inférieure : il se tient là où le modèle attend les trois
autres maisons.
\end{chiffres}

\begin{resultat}[Conséquence sur l'entraînement]
Le coût, et donc la correction qui en découle, dépend principalement d'un petit
nombre d'élèves atypiques. Déplacer la frontière pour reclasser ces quelques points
fait baisser $J$ davantage qu'augmenter la probabilité accordée aux mille cinq
cents autres. De là vient sa sensibilité aux valeurs aberrantes, ce qui impose
d'examiner les élèves mal classés et pas seulement la moyenne.
\end{resultat}

\section{Comptage et coût}

Savoir si une frontière est bien placée admet deux réponses. Compter les élèves
du bon côté est la plus directe : le résultat est un nombre entier, qui se lit
sans convention. Faire la moyenne des coûts individuels donne un nombre réel,
dont la valeur ne s'interprète pas seule. C'est pourtant la seconde qui sert à
l'entraînement, et la raison tient à la manière dont chacune réagit lorsque la
frontière se déplace.

Le comptage change par sauts. Tant que la frontière se déplace sans que personne
ne la traverse, il garde en effet la même valeur, et il saute d'une unité à
l'instant où un élève passe de l'autre côté. Sur tout l'intervalle qui sépare
deux sauts, il rend la même réponse pour des frontières différentes, et un
calcul qui chercherait à l'améliorer recevrait partout la même information.

Le coût répond à tout déplacement. Chaque élève y contribue par la probabilité
que le modèle lui accorde, et cette probabilité change dès que la frontière
bouge, fût-ce d'un millième d'écart-type. La moyenne bouge elle aussi, et le
sens dans lequel elle bouge désigne le sens du déplacement à faire.

Trois jeux de coefficients, mesurés sur les mêmes $1\,600$ élèves, donnent à voir
cette différence.

\begin{table}[htbp]
\centering
\begin{tabular}{lccc}
\toprule
coefficients & frontière & élèves bien classés & coût $J$ \\
\midrule
$w$ au $50^{\text{e}}$ tour    & \multicolumn{1}{c}{---}        & $1\,563$ & $0{,}1438$ \\
$w$ au $400^{\text{e}}$ tour   & voisine de la précédente       & $1\,564$ & $0{,}1081$ \\
$1{,}5\,w$ au $400^{\text{e}}$ & identique à la précédente      & $1\,564$ & $0{,}1155$ \\
\bottomrule
\end{tabular}
\caption{Les deux premières lignes séparent presque les mêmes élèves ; la
troisième multiplie les coefficients par $1{,}5$, ce qui laisse la frontière
exactement où elle est.}
\label{tab:comparaison}
\end{table}

Les deux premières lignes portent des frontières voisines, que le comptage
sépare d'une unité et le coût d'un quart de sa valeur : trois cent cinquante
tours de travail apparaissent dans une seule des deux mesures. La troisième garde
la frontière exactement où elle était, en multipliant les trois coefficients par
un même facteur --- $w$ et $\lambda w$ s'annulent aux mêmes points. Les mêmes
élèves sont du bon côté, le comptage rend le même nombre, et le coût distingue
pourtant les deux modèles : l'écartement des probabilités profite peu aux élèves
déjà bien classés et pénalise lourdement les trente-six autres.

\begin{figure}[htbp]
\centering
\begin{tikzpicture}
\begin{axis}[width=.80\textwidth,height=6.4cm,
  axis lines=left,axis line style={ink!40},
  grid=major,grid style={gridgray!30},
  tick label style={font=\small},label style={font=\small},
  xlabel={biais $w_0$},ylabel={coût $J$},
  xmin=-6.5,xmax=-2.6,ymin=0.10,ymax=0.23,
  axis y line*=left]
  \addplot[accent,very thick,no markers]
    table[x=w0,y=cost] {data/comparaison.dat};
\end{axis}
\begin{axis}[width=.80\textwidth,height=6.4cm,
  axis lines=right,axis line style={warning!50},
  tick label style={font=\small,text=warning},
  label style={font=\small,text=warning},
  ylabel={élèves mal classés},
  xmin=-6.5,xmax=-2.6,ymin=25,ymax=150,
  axis x line=none,axis y line*=right]
  \addplot[warning,thick,no markers,const plot]
    table[x=w0,y=wrong] {data/comparaison.dat};
\end{axis}
\end{tikzpicture}
\caption{Les deux mesures le long d'un déplacement du biais, les deux autres
coefficients restant à leur valeur finale. Le nombre de mal classés progresse par
paliers --- $68$ valeurs distinctes sur les $161$ positions essayées --- et reste
plat sur des intervalles entiers : aucune correction ne s'en déduit. Le coût
varie à chaque déplacement, ce qui permet d'en tirer une direction. Les deux minima tombent à des endroits différents : le coût est le plus bas en $w_0=-4{,}80$, les erreurs les moins nombreuses en $w_0=-5{,}68$.}
\label{fig:deux-mesures}
\end{figure}
\FloatBarrier

C'est cette sensibilité que le chapitre suivant exploite : une quantité qui
répond à tout déplacement, si petit soit-il, se dérive, et sa dérivée indique
dans quel sens déplacer les coefficients.
